\newpage
\section{Incidência da empenagem horizontal}

A incidência da empenagem horizontal de cada configuração foi ajustada para
anular a deflexão de profundor na condição de cruzeiro. Em cada arquivo
fixou-se a meta de sustentação do ponto de projeto ($M = 0{,}85$,
$C_L = 0{,}5053$) e a restrição de momento de arfagem nulo pelo profundor.
Como a resposta $\delta_e(i_t)$ é linear no método de malha de vórtices, o
zero foi resolvido com duas avaliações e uma rodada de verificação,
resultando em deflexões residuais abaixo de $0{,}01^\circ$, uma ordem de
grandeza melhor que a tolerância de $0{,}1^\circ$ sugerida no roteiro. A
Tabela~\ref{tab:trimagem} resume os resultados para a asa final, com a
torção otimizada da Seção~\ref{sec:torcao}.

\begin{table}[H]
\centering
\caption{Incidência da empenagem que anula a deflexão de profundor em cruzeiro.}\label{tab:trimagem}
\renewcommand{\arraystretch}{1.25}
\begin{tabular}{lccc}
\toprule
Configuração & $i_t$ [graus] & $\delta_e$ residual [graus] & $\alpha$ [graus] \\
\midrule
CG dianteiro (\texttt{fwd.avl}) & $-4{,}48$ & $-0{,}001$ & 5,38 \\
CG traseiro (\texttt{aft.avl}) & $-1{,}98$ & $+0{,}0004$ & 5,07 \\
\bottomrule
\end{tabular}
\end{table}

As incidências resultantes ficaram próximas da faixa de um transporte
convencional, que tipicamente opera com estabilizador entre $-3^\circ$ e
$+1^\circ$, com a torção da asa aliviando cerca de $4{,}0^\circ$ em
relação à asa sem torção, um dos ganhos da otimização da Seção 6. A
magnitude residual, ainda algo alta no CG dianteiro, tem dois fatores.
Primeiro, a fuselagem não é modelada no AVL, o que desloca o ponto neutro
para trás e infla a margem estática, exigindo mais força de compensação da
empenagem. Segundo, os perfis supercríticos otimizados no Lab 03 têm
momento de arfagem fortemente negativo, que precisa ser equilibrado pela
cauda. A validação do ajuste aparece nas polares da
Seção~\ref{sec:polares}: nas curvas de $C_L \times \delta_e$ os casos
compensados cruzam $\delta_e = 0$ exatamente no $C_L$ de projeto.
